\chapter{Nutzung von KI-Werkzeugen in deiner Abschlussarbeit}
\label{chap:ki-werkzeuge}

KI-Werkzeuge (z.\,B. ChatGPT, Claude, Copilot, Claude Code, Cursor, DeepL, Perplexity, Midjourney) gehören heute zum Handwerkszeug wissenschaftlichen Arbeitens und dürfen in deiner Abschlussarbeit selbstverständlich eingesetzt werden. Ich ermutige dich ausdrücklich, diese Tools produktiv zu nutzen -- sei es zur Ideenfindung, als Sparringspartner beim Strukturieren, für Grammatik- und Stilkorrekturen, zur Literaturrecherche, für Code-Unterstützung, beim Einsatz von Coding-Agents (z.\,B. zur Umsetzung ganzer Module, für Refactoring, Debugging oder das Aufsetzen von Tests), zur Visualisierung oder für Übersetzungen.

Zwei Punkte sind mir dabei wichtig.

\section{Du verantwortest das Ergebnis}

Alle Aussagen, Argumente, Codezeilen und Abbildungen in deiner Arbeit musst du inhaltlich verstehen, prüfen und vertreten können. KI-Ausgaben sind nie ungeprüft zu übernehmen -- insbesondere Quellenangaben, Zahlen und fachliche Behauptungen verifiziere bitte eigenständig. Für Code, der von Agents generiert oder umgebaut wurde, gilt dasselbe: Du musst ihn lesen, verstehen, testen und bei der Verteidigung erklären können.

\section{Transparenz statt Bürokratie}
\label{sec:ki-transparenz}

Füge deiner Arbeit einen kurzen Anhang \enquote{Einsatz von KI-Werkzeugen} bei (siehe Vorlage im Anhang dieser Arbeit). Eine knappe Tabelle oder ein paar Sätze genügen, in denen du festhältst:

\begin{itemize}
    \item[a)] welche Tools du verwendet hast,
    \item[b)] wofür (z.\,B. Brainstorming, Sprachglättung, Code-Refactoring, agentische Code-Generierung, Bildgenerierung, Literaturrecherche),
    \item[c)] an welchen Stellen der Einsatz substanziell war -- Kapitel- oder Abschnittsangaben reichen, bei Code z.\,B. die betroffenen Module oder Verzeichnisse; du musst nicht jede Formulierung oder Zeile markieren.
\end{itemize}

Bitte ergänze im Anhang folgenden Satz:

\begin{quote}
\enquote{Ich habe die eingesetzten KI-Werkzeuge im beschriebenen Umfang genutzt und verantworte Auswahl, Übernahme und inhaltliche Korrektheit der Ergebnisse vollumfänglich selbst.}
\end{quote}

\section{Was du nicht einzeln ausweisen musst}

Rechtschreib- und Grammatikprüfung, Autovervollständigung oder kurze Wortvorschläge -- auch bei Code-Editoren, die zeilenweise vervollständigen. Alles, was darüber hinaus die Substanz der Arbeit berührt, gehört in die Übersicht im Anhang.

\section{Was nicht zulässig ist}

Nicht zulässig ist lediglich, KI-generierte Inhalte als ausschließlich eigene Leistung auszugeben und den Einsatz zu verschweigen. Solange du offen dokumentierst, wie KI dich unterstützt hat, ist die Nutzung willkommen -- und oft sogar wünschenswert.
